\documentclass[12pt,a4paper]{article}
\usepackage[utf8]{inputenc}
\usepackage[T1]{fontenc}
\usepackage[romanian]{babel}
\usepackage[margin=2.5cm]{geometry}
\usepackage{graphicx}
\usepackage{listings}
\usepackage{xcolor}
\usepackage{caption}

\newcommand{\studentgroup}{FAF-242}
\definecolor{termbg}{rgb}{0.96,0.96,0.96}
\definecolor{codegray}{rgb}{0.5,0.5,0.5}

\lstset{
  backgroundcolor=\color{termbg},
  basicstyle=\ttfamily\footnotesize,
  breaklines=true,
  captionpos=b,
  frame=single,
  keepspaces=true,
  numbers=left,
  numbersep=6pt,
  numberstyle=\tiny\color{codegray},
  showstringspaces=false,
  tabsize=2
}

\begin{document}

% =========================================================================
% TITLE PAGE
% =========================================================================
\begin{titlepage}
    \centering
    \IfFileExists{utm_logo(1).png}{%
        \includegraphics[width=0.25\textwidth]{utm_logo(1).png}\\[0.6cm]%
    }{%
        \vspace*{2cm}%
    }
    {\large \textbf{UNIVERSITATEA TEHNICĂ A MOLDOVEI}}\\[0.2cm]
    {\large Facultatea Calculatoare, Informatică și Microelectronică}\\[0.2cm]
    {\large Departamentul Ingineria Software și Automatică}\\[2.5cm]

    {\Large \textbf{LUCRAREA DE LABORATOR NR. 9 \& 10}}\\[0.4cm]
    {\LARGE \textbf{Watching System Calls cu strace (Linux) \&\\ Add a System Call: trace (xv6)}}\\[0.4cm]
    {\large Disciplina: Sisteme de Operare}\\[2.5cm]

    \begin{minipage}{0.45\textwidth}
        \flushleft
        \textbf{Done by:}\\
        st. gr. \studentgroup\\
        Daniel-Ioan \textsc{Strunga}
    \end{minipage}
    \hfill
    \begin{minipage}{0.50\textwidth}
        \flushright
        \textbf{Checked by:}\\
        asist. univ. Patricia \textsc{Reițman}\\
        asist. univ. Marius \textsc{Dumitrașcu}
    \end{minipage}

    \vfill
    {\large Chișinău --- 2026}
\end{titlepage}

\newpage
\tableofcontents
\newpage

\section{Scopul lucrării}
Această lucrare combinată are un dublu obiectiv. Pe de o parte (Lab 9 - Track A), se urmărește observarea apelurilor de sistem declanșate de programe reale în Linux folosind utilitarul \texttt{strace} și exersarea apelurilor POSIX fundamentale de I/O (\texttt{open}, \texttt{read}, \texttt{write}, \texttt{close}) omițând bibliotecile de nivel înalt. Pe de altă parte (Lab 10 - Track B), se vizează extinderea nucleului xv6 prin implementarea unui apel de sistem complet nou (\texttt{trace}) capabil să monitorizeze execuția altor syscall-uri pe baza unor măști de biți moștenite între procese.

% =========================================================================
% LAB 9: LINUX
% =========================================================================
\section{Laboratorul 9: Watching System Calls (Linux)}

\subsection{Codul sursă POSIX I/O (\texttt{syscalls.c})}
Acest program demonstrează scrierea și citirea folosind exclusiv apeluri de sistem native, fără interfața \texttt{FILE*} din \texttt{stdio.h}.
\begin{lstlisting}[language=C]
#include <fcntl.h>
#include <unistd.h>
#include <stdio.h>

int main(void) {
    // Creare si scriere
    int fd = open("note.txt", O_WRONLY | O_CREAT | O_TRUNC, 0644);
    if (fd < 0) { perror("open write"); return 1; }
    printf("open returned descriptor %d\n", fd);
    write(fd, "hello, kernel\n", 14);
    printf("write sent 14 bytes\n");
    close(fd);

    // Redeschidere si citire
    fd = open("note.txt", O_RDONLY);
    if (fd < 0) { perror("open read"); return 1; }
    char buf[16];
    read(fd, buf, 14);
    buf[14] = '\0';
    printf("read got 14 bytes: %s", buf);
    close(fd);
    
    return 0;
}
\end{lstlisting}

\subsection{Ieșirile din terminal (strace)}
\begin{figure}[h!]
\begin{lstlisting}[language=bash]
$ strace -c ls -l >/dev/null
% time     seconds  usecs/call     calls    errors syscall
------ ----------- ----------- --------- --------- ----------------
 17.63    0.000213          12        17           mmap
 13.41    0.000162          10        16           close
  8.36    0.000101           9        11           read
  7.78    0.000094           9        10           openat
  0.83    0.000010          10         1           write
------ ----------- ----------- --------- --------- ----------------
100.00    0.001208                   117        13 total
\end{lstlisting}
\caption{Figure 9.1 -- Statistica syscall-urilor pentru ls -l pe arhitectura x86\_64 (117 apeluri totale)}
\end{figure}

\begin{figure}[h!]
\begin{lstlisting}[language=bash]
$ strace -e trace=openat,write,read,close ./syscalls
openat(AT_FDCWD, "/etc/ld.so.cache", O_RDONLY|O_CLOEXEC) = 3
close(3)                                = 0
openat(AT_FDCWD, "/lib/x86_64-linux-gnu/libc.so.6", O_RDONLY|O_CLOEXEC) = 3
read(3, "\177ELF\2\1\1\3\0\0\0\0\0\0\0\0\3\0>\0\1\0\0\0\220\243\2\0\0\0\0\0"..., 832) = 832
close(3)                                = 0
openat(AT_FDCWD, "note.txt", O_WRONLY|O_CREAT|O_TRUNC, 0644) = 3
write(1, "open returned descriptor 3\n", 27) = 27
write(3, "hello, kernel\n", 14)         = 14
write(1, "write sent 14 bytes\n", 20)   = 20
close(3)                                = 0
openat(AT_FDCWD, "note.txt", O_RDONLY)  = 3
read(3, "hello, kernel\n", 63)          = 14
write(1, "read got 14 bytes: hello, kernel"..., 33) = 33
close(3)                                = 0
+++ exited with 0 +++
\end{lstlisting}
\caption{Figure 9.2 -- Urmărirea programului C propriu. Se observă cum kernel-ul refolosește descriptorul 3.}
\end{figure}

\subsection{Răspunsuri (Observe Lab 9)}
\begin{enumerate}
    \item \textbf{Cum afișează \texttt{ls} pe ecran?} Comanda \texttt{ls} scrie conținutul pe ecran apelând funcția de sistem \texttt{write}, direcționând datele către descriptorul de fișier \texttt{1} (\texttt{stdout}).
    \item \textbf{De ce fișierul primește descriptorul 3?} La pornirea oricărui proces, kernel-ul asociază implicit primii trei descriptori: 0 (\texttt{stdin}), 1 (\texttt{stdout}) și 2 (\texttt{stderr}). Deoarece regula de alocare a sistemului de operare este să ofere cel mai mic număr liber disponibil, primul fișier deschis prin \texttt{openat} primește \textbf{3}. După execuția lui \texttt{close(3)}, numărul redevine disponibil, fiind refolosit la a doua apelare \texttt{openat}.
    \item \textbf{De ce obținem rezultate diferite fără \texttt{>/dev/null}?} Când nu redirecționăm output-ul (\texttt{>/dev/null}), textul printat de propriul nostru program, de exemplu linia \textit{„write sent 14 bytes\textbackslash n”}, va fi interceptat și trimis către terminal. Deoarece linia conține exact cuvântul \texttt{write} la început, comanda \texttt{grep '\textasciicircum write'} face un match fals pe output-ul propriu-zis al programului, numărând eronat 3 apariții în loc de 2.
\end{enumerate}

\newpage

% =========================================================================
% LAB 10: XV6
% =========================================================================
\section{Laboratorul 10: Add a System Call: trace (xv6)}

\subsection{Codul sursă user-space (\texttt{user/trace.c})}
\begin{lstlisting}[language=C]
#include "kernel/param.h"
#include "kernel/types.h"
#include "kernel/stat.h"
#include "user/user.h"

int main(int argc, char *argv[]) {
  if (argc < 3) {
    fprintf(2, "Usage: %s mask command\n", argv[0]);
    exit(1);
  }
  
  if (trace(atoi(argv[1])) < 0) {
    fprintf(2, "%s: trace failed\n", argv[0]);
    exit(1);
  }
  
  exec(argv[2], &argv[2]);
  exit(0);
}
\end{lstlisting}

\subsection{Ieșirile din terminal (xv6 emulator)}
\begin{figure}[h!]
\begin{lstlisting}[language=bash]
$ trace 32 grep hello README
3: syscall read -> 1023
3: syscall read -> 965
3: syscall read -> 453
3: syscall read -> 0
\end{lstlisting}
\caption{Figure 10.1 -- Izolarea apelurilor de sistem 'read' folosind masca pe biți}
\end{figure}

\begin{figure}[h!]
\begin{lstlisting}[language=bash]
$ trace 2 usertests forkforkfork
usertests starting
test forkforkfork: 5: syscall fork -> 6
6: syscall fork -> 7
7: syscall fork -> 8
8: syscall fork -> 9
...
26: syscall fork -> 42
...
9: syscall fork -> 66
9: syscall fork -> -1
OK
ALL TESTS PASSED
\end{lstlisting}
\caption{Figure 10.2 -- Verificarea mecanismului de moștenire a măștii în timpul creării de procese multiple}
\end{figure}

\begin{figure}[h!]
\begin{lstlisting}[language=bash]
$ trace 2147483647 grep hello README
4: syscall trace -> 0
4: syscall exec -> 3
4: syscall open -> 3
4: syscall read -> 1023
4: syscall read -> 965
4: syscall read -> 453
4: syscall read -> 0
4: syscall close -> 0
\end{lstlisting}
\caption{Figure 10.3 -- Monitorizarea tuturor apelurilor cu masca maximă (toți biții setați pe 1)}
\end{figure}

\subsection{Răspunsuri (Observe Lab 10)}
\begin{enumerate}
    \item \textbf{De ce 32 izolează apelul \texttt{read}?} Valoarea $32$ reprezintă în sistem zecimal numărul $2^5$. Apelul \texttt{SYS\_read} este definit în sistem cu ID-ul $5$. Shiftarea la stânga (\texttt{1 << 5}) setează exact bitul corespunzător acestui apel, ignorând restul operațiunilor.
    \item \textbf{Cum se moștenește masca prin \texttt{fork}?} Atunci când un proces se clonează, kernel-ul rulează funcția \texttt{kfork()}. Acolo am adăugat manual atribuirea \texttt{np->trace\_mask = p->trace\_mask;}, ceea ce face ca noul proces (copilul) să copieze exact setările de monitorizare ale părintelui său. Această stare este reținută în structura \texttt{proc} din kernel.
    \item \textbf{De ce \texttt{echo hi} generează două apeluri \texttt{write}?} Programul \texttt{echo} este scris să afișeze mai întâi șirul de caractere trimis ca argument (\texttt{hi}), consumând un apel \texttt{write}. Ulterior, efectuează un al doilea apel separat pentru a tipări terminatorul de rând (newline, \texttt{\textbackslash n}).
    \item \textbf{De ce se loghează \texttt{printk} după apelul sistem?} Am inserat instrucțiunea de logare \texttt{printk} în \texttt{syscall.c} \textit{după} execuția rutinei efective a apelului. Motivul este că aveam nevoie să preluăm și să printăm valoarea returnată de acel apel (ex: numărul de octeți citiți), care este salvată de kernel în registrul \texttt{p->trapframe->a0} abia la finalizarea execuției.
    \item \textbf{Cum formăm masca pentru \texttt{open} și \texttt{close}?} \texttt{SYS\_open} are ID-ul 15 ($2^{15} = 32768$), iar \texttt{SYS\_close} are ID-ul 21 ($2^{21} = 2097152$). Adunând aceste două valori obținem masca compusă \textbf{2129920}.
\end{enumerate}

\newpage

% =========================================================================
% CONCLUZII FINALE
% =========================================================================
\section{Concluzii (Comparare Linux - xv6)}
Prin desfășurarea laboratoarelor 9 și 10 am studiat interfața nativă de comunicare dintre aplicațiile din user-space și kernel, implementată printr-un perete de securitate bazat pe apeluri de sistem (syscall-uri).

\textbf{Diferența de mediu (strace vs trace):}
Pe sistemul de operare complet echipat Linux (Lab 9), ne-am bazat pe uneltele de audit existente (precum \texttt{strace}). Am observat cantitatea uriașă de muncă depusă în fundal doar pentru a rula \texttt{ls -l} (117 apeluri pe x86\_64), mare parte din ele implicând încărcarea bibliotecilor partajate (fișiere \texttt{.so}) dinamic. \texttt{libc} joacă un rol major de abstractizare și buffering pentru a reduce suprasolicitarea kernel-ului.

Pe sistemul educațional xv6 (Lab 10), arhitectura fiind minimalistă, a fost nevoie să intervenim direct în fișierele C ale kernelului (\texttt{sysproc.c}, \texttt{syscall.c}) pentru a dezvolta manual un subsistem de audit, capabil să prindă tranzițiile în momentul în care procesorul schimbă modul de privilegiu (trap). Implementarea pe bază de mască bitwise s-a dovedit a fi un mecanism extrem de performant ca timp de execuție ($O(1)$) și eficient la nivel de alocare de memorie pentru definirea stării unui proces monitorizat.

Ambele laboratoare au revalidat structura uniformă a descriptorilor de fișiere POSIX: zero, unu și doi sunt alocați invariabil fluxurilor standard de la inițializarea procesului, confirmând axioma fundamentală a arhitecturilor derivate din UNIX.

\end{document}